\documentclass[10pt,a4paper]{amsart}
\usepackage[margin=19mm]{geometry}
\usepackage{amsmath,amssymb,mathtools}
\usepackage[scr=rsfs,cal=euler]{mathalfa}
\usepackage{xfrac}
\usepackage[unicode,hidelinks,bookmarks=false]{hyperref}
\newcommand{\ie}{i.e.\ }
\newcommand{\cf}{cf.\ }
\newcommand{\resp}{respectively\ }
\newcommand{\Subsection}{\subsection}
\providecommand{\nobreakdash}{\nobreak\hbox{-}}
\setlength{\emergencystretch}{2em}
\multlinegap0pt
\usepackage{amssymb}
\newtheorem{lemma}{Lemma}
\newcommand{\C}{{\mat C}}
\providecommand{\Diag}{\operatorname{diag}}% diagonal
\newcommand{\Om}{\Omega}
\providecommand{\Rp}{\operatorname{Re}}% Real part
\newcommand{\ba}{\backslash}
\newcommand{\bs}{\boldsymbol}
\newcommand{\dive}{{\rm div}}
\newcommand{\mat}{\mathbb}
\newcommand{\ov}{\overline}
\newcommand{\pa}{\partial}
\newcommand{\sig}{\sigma}
\newcommand{\wid}{\widetilde}
\title{\bf Convergence of the PML method for scattering problems in poroelastic media}
\author{Qianyuan Yin, Changkun Wei, Bo Zhang}
\date{Source: arXiv:2606.30392v1 (CC BY 4.0)}
\begin{document}
\maketitle

\noindent\emph{Source and licence.} \bf Convergence of the PML method for scattering problems in poroelastic media. Version arXiv:2606.30392v1 (CC BY 4.0), distributed by arXiv under the Creative Commons Attribution 4.0 International licence, http://creativecommons.org/licenses/by/4.0/, as recorded in the arXiv licence metadata for this exact version and shown on the abstract page https://arxiv.org/abs/2606.30392v1. Full licence notice: \texttt{licenses/arxiv-2606.30392-CC-BY-4.0.txt}.\par\smallskip

\noindent\textit{Editorial scope.} Counted unit: the article's lemma PMLHD (source0.tex lines 1153--1163). The article's complete original proof of it is delivered with it (source0.tex lines 1164--1220, one \texttt{proof} environment of 57 lines). No mathematics is written, repaired or re-derived: the statement and the proof are the article's own lines, in the article's own notation, and no retained line is substituted or re-flowed.  Retained uncounted as the source-local context the proof needs: lines 1000--1002; lines 1070--1073; lines 1075--1077; lines 1079--1083.  Nothing was omitted from any retained span, so the package carries no omission record.  No retained line contains a citation, so the wrapper carries no bibliography and no bibliography entry.  No retained line contains a figure environment, an image-inclusion command or drawing code, so no image is delivered and the package declares no asset dependency.  Editorial formatting only, in addition: an AMS \texttt{amsart} preamble that restates the article's own declarations and macros used by the retained lines, namely the environments \texttt{lemma} on separate counters, together with the standard packages those lines need.  The retained environments print in this document as Lemma 1 in order of appearance, whereas the article prints the same items with the numbering of its own full document.  The complete native source of the article as distributed in the arXiv e-print for this exact version is bundled unchanged as \texttt{sources/204016/source0.tex}.
\begin{equation}\label{stretched_r}
\wid{r}(r)=\int_0^{r}s(\tau)d\tau \triangleq r\beta(r) = r[1+{\bf i}\sig_m(r)],
\end{equation}

\begin{equation}\label{relation_oper}
\wid{\nabla}f=P_1^{-1}\nabla f,\;\wid{\dive}\mathbf{u}=\frac{1}{\beta^2s}
\dive (P_2 \mathbf{u}),\;\wid{\nabla}\times\mathbf{u}=P_2^{-1}\nabla \times P_1\mathbf{u},
\end{equation}

\begin{equation}\label{P1P2}
P_1=\Diag\{s,\beta,\beta\},\quad P_2=\Diag\{\beta^2,s\beta,s\beta\}.
\end{equation}

\begin{equation}\label{stretched_laplace}
\wid{\Delta}\mathbf{u}=\sum_{i=1}^3(\wid{\Delta}u_i)\bs{e}_i
=\sum_{i=1}^3\wid{\dive}(\wid{\nabla}u_i)\bs{e}_i=- \wid{\nabla}\times\wid{\nabla}\times u
+\wid{\nabla} \wid{\dive}\mathbf{u}.
\end{equation}

\begin{lemma}[PML Helmholtz decomposition]\label{PMLHD}
Let $D=B_{r_2}\ba\ov{\Om}$ or $D=B_{r_2}\ba\ov{B_{r_0}}$ be the truncated PML domain or the PML layer.
Define $V(D)=\{\bs{v}\in H^1(D)^3:\bs{n} \cdot\bs{v} =0 \;{\rm on}\;\pa D\}
$. The following two decompositions hold:
\begin{enumerate}
\item[(1)] For any $\bs{U}\in H_0^1(D)^3$, there exists $\theta_1 \in H^2(D)\cap H_0^1(D)$ and $\Phi_1\in H^1(D)^3$ 
such that $\bs{U}= \wid{\nabla}\theta_1 +\Phi_1$ and $\wid{\dive}\Phi_1=0$.
\item[(2)] For any $\bs{W}\in V(D)$, there exists $\theta_2 \in H^2(D)/\mathbb{C}$ and $\Phi_2\in V(D)$ such that 
$\bs{W}= \wid{\nabla}\theta_2 +\Phi_2$, $\wid{\dive}\Phi_2=0$ and $\bs{n}\cdot\wid{\nabla}\theta_2=0$ on $\pa D$.
\end{enumerate}
\end{lemma}

\begin{proof}
(1) For any fixed $\bs{U}\in H_0^1(D)^3$, consider the Dirichlet problem
\begin{equation}\label{eq:dirichlet}
\wid{\Delta}\theta = \wid{\dive}\bs{U} \quad\text{in } D, \qquad \theta = 0 \quad\text{on } \pa D .
\end{equation}
From the relations of the stretched operators \eqref{relation_oper}-\eqref{stretched_laplace}, 
we have the identity
\begin{equation}
\wid{\Delta}\theta = \frac{1}{s\beta^2}\dive\bigl(P_2P_1^{-1}\nabla\theta\bigr).
\end{equation}
This leads to the weak formulation of \eqref{eq:dirichlet}: to find $\theta\in H_0^1(D)$ such that
\begin{equation*}
a(\theta,\psi):=\int_D P_2P_1^{-1}\nabla\theta\cdot\nabla\overline{\psi}\,dx
= -\int_D s\beta^2\,(\wid{\dive}\bs{U})\,\overline{\psi}\,dx
\quad \forall \psi\;\in H_0^1(D).
\end{equation*}
It follows from the definitions of $P_1$ and $P_2$ (see \eqref{P1P2}) that
\begin{equation}\label{eq:variation}
\begin{aligned}
\Rp a(\psi,\psi)&=\Rp\int_{D}P_2P_1^{-1}\nabla \psi\cdot\nabla\overline{\psi}dx \\
&= \Rp\int_{D}\Diag\{\beta^2s^{-1},s,s\}\nabla \psi\cdot\nabla\overline{\psi}dx \\
&=\Rp\int_{D}\left[\beta^2s^{-1}\vert (\nabla \psi)_r \vert^2+s\vert (\nabla \psi)_\theta\vert^2
+s\vert (\nabla \psi)_\phi \vert^2\right]dx.
\end{aligned}
\end{equation}
From the definition of $\beta$ in \eqref{stretched_r} and the fact that $\sig_m\le\sig$ for $r\ge r_0$, we obtain
\begin{equation}\label{bs_lower_bound}
\Rp(\beta^2s^{-1}) = \Rp\frac{(1+{\bf i}\sig_m)^2}{1+{\bf i}\sig} = \frac{1-\sig_m^2+2\sig\sig_m}{1+\sig^2} 
\geq \frac{1}{1+\Vert\sig\Vert_{L^{\infty}(B_{r_2})}^2}.
\end{equation}
This, combined with \eqref{eq:variation} yields
\begin{equation}\label{lowerbound_a}
\Rp a(\theta,\theta) \geq
\frac{\lVert \nabla \theta \rVert_{{L^2(D)^3}}^2}{1+\Vert\sig\Vert_{L^{\infty}(B_{r_2})}^2}.
\end{equation}
By the Poincar\'e inequality and elliptic regularity theory, problem \eqref{eq:dirichlet} 
admits a unique solution $\theta_1\in H^2(D)\cap H_0^1(D)$. Setting $\Phi_1 = \bs{U} - 
\wid{\nabla}\theta_1$ gives the desired decomposition.

(2) For any fixed $\bs{W}\in V(D)$, consider the Neumann problem
\begin{equation}\label{eq:neumann}
\wid{\Delta}\theta=\wid{\dive}\bs{W} \quad\text{in } D, \qquad \bs{n}\cdot \wid{\nabla}\theta=0 
\quad\text{on } \pa D.
\end{equation}
The compatibility conditions are satisfied since
\begin{equation*}
\begin{aligned}
\int_{D}\dive(P_2P_1^{-1}\nabla\theta)dx &=\int_{\pa D}\bs{n}\cdot P_2P_1^{-1}\nabla\theta ds(x)= 
\int_{\pa D}\frac{\beta^2}{s}\bs{n}\cdot \nabla\theta ds(x)=0\\
\int_{D}\dive(P_2\bs{W})dx&=\int_{\pa D}\bs{n}\cdot P_2\bs{W}ds(x)= \int_{\pa D}\beta^2\bs{n}\cdot\bs{W}ds(x)=0.
\end{aligned}
\end{equation*}
Proceeding as in the first part, the Lax-Milgram theorem and elliptic regularity theory implies the 
 existence of the unique solution
$\theta_2 \in H^2(D)/\C$ to the problem \eqref{eq:neumann}. Defining $\Phi_2 = \bs{W} - \wid{\nabla}\theta_2$ 
completes the proof.
\end{proof}

\end{document}
